% Lösungen Einheit 3 von 5 – Additionsverfahren (zusammenbau v0.1)
\einheitenkopf{Einheit 3 von 5 · Additionsverfahren}

\erg{19}{a) Gegenzahlen \quad b) $x = 3$, $y = 5$; Lösung $(3|5)$ \quad c) $x = 5$, $y = 2$; Lösung $(5|2)$ \quad d) $x = 3$, $y = 1$; Lösung $(3|1)$ \quad e) $x = 5$, $y = 2$; Lösung $(5|2)$ \quad f) $x = -1$, $y = 4$; Lösung $(-1|4)$ \quad g) $x = 4$, $y = 2$ \quad h) $2 \cdot$ I $-$ II: $0 = 3$ (fA): keine Lösung \quad i) Gleichsetzen: Beide Gleichungen sind schon nach $y$ aufgelöst. \quad j) $2 \cdot$ I $+$ II: $8x = 40$, $x = 5$, $y = 2$; Probe: $15 + 8 = 23$ (wA), $10 - 16 = -6$ (wA). Addition, weil keine Variable steht allein, nach dem Verdoppeln von I stehen bei $y$ Gegenzahlen.}
\erg{20}{a) Gegenzahlen \quad b) I $+$ II: $0 = 5$ (fA): keine Lösung \quad c) $9$ (II ist das Dreifache von I) \quad d) Zum Beispiel $a = -6$, $b = 0$: Die linke Seite von II ist das $(-3)$-Fache von I, die rechte nicht. $3 \cdot$ I $+$ II ergibt $0 = 12$ (fA), also keine Lösung (jedes $b \ne -12$ geht).}
\erg{21}{a) $x = 4$, $y = 3$; Lösung $(4|3)$}
\erg{22}{a) Beim Verdoppeln muss auch die rechte Seite mal $2$: $2x + 4y = 14$; dann $5x = 15$, $x = 3$, $y = 2$. \quad b) $3y - (-y) = 3y + y$: richtig ist $4y = 16$, also $y = 4$, $x = 1$. \quad c) Mit $a = 2$ ist die linke Seite von II kein Vielfaches von I, das System bleibt eindeutig lösbar. Richtig: $a = -2$, dann ergibt $2 \cdot$ I $+$ II $0 = 5$ (fA). \quad d) Für die gesuchte Lösung sind in I und in II beide Seiten gleich groß. Addiert man Gleiches zu Gleichem, sind auch die Summen gleich – die Lösung erfüllt die neue Gleichung. \quad e) Hätte das System eine Lösung, müsste sie auch die addierte Gleichung erfüllen. $0 = 4$ ist aber für kein Zahlenpaar wahr – also gibt es keine Lösung.}
